\section{Problem Definition}
\label{sec:goals_tasks}
Designing grippers for a given set of workpieces is still largely a manual, experience-driven task, in which geometric parameters such as suction-cup
placement or finger contours are chosen iteratively and validated in simulation or on hardware\cite{Liu2022Simulation}. Because grasp quality can only
be assessed through a simulation or geometric evaluation routine, no analytical gradients or closed-form objective expressions are available, and each
evaluation is comparatively expensive\cite{Santoni2023Comparison}. Moreover, a single gripper usually has to handle several workpieces at once, so the
design criteria are inherently conflicting and cannot be reduced to a single scalar quality measure without losing information about the available
trade-offs. These properties of expensive, derivative-free, multi-objective evaluations place gripper design squarely in the domain of black-box
multi-objective optimization\cite{Mueller2017SOCEMO}. Which class of optimization algorithm is best suited for such problems, however, remains an open
question, as performance is known to depend strongly on the number of objectives, the shape of the objective landscape, and the available evaluation
budget\cite{Santoni2023Comparison, Raponi2023Optimizing}.

The primary goal of this thesis is to apply and comparatively evaluate optimization algorithms spanning evolutionary multi-objective methods and
Bayesian optimization for solving gripper configuration problems formulated as black-box multi-objective optimization
tasks\cite{Datta2011Optimizing,Belakaria2020Uncertainty-Aware}. The two gripper types, a multi-suction-cup gripper and a two-finger gripper, serve as
well-defined benchmark problems on which the algorithms are systematically compared with respect to solution quality, convergence behaviour, and
robustness.

To provide a rigorous and reproducible basis for this comparison, two gripper configuration problems are formulated as well-defined multi-objective
optimization tasks. In both cases, a diverse set of workpieces is incorporated directly into the problem formulation by treating each workpiece as a
separate objective. This design choice ensures that the optimized gripper configurations are not tailored to a single
geometry, but are instead robust across a range of workpiece shapes. It also increases the dimensionality of the objective space as the number of
workpieces grows, allowing the benchmarking study to probe algorithm behaviour under varying levels of problem complexity\cite{Guo2022Evolutionary}.

For the \textbf{multi-suction-cup gripper}, the objectives consist of the geometric stability of the suction base together with the grasp diversity
evaluated independently for each workpiece. These objectives are inherently conflicting: a larger suction base improves stability but increases the
risk of collisions with the workpiece, reducing the number of feasible grasps.

For the \textbf{two-finger gripper}, a fixed number of grasp points is defined for each workpiece, and the Objectives are computed from the conformance of the
finger contour to the corresponding workpiece section at each grasp point. A finger contour that scores well at all grasp points simultaneously must
therefore generalize across diverse geometric profiles, which makes the problem increasingly challenging as the number of grasp points grows.

Together, these two problems provide a structured experimental basis for comparing optimization algorithms across varying objective space
dimensionalities and objective landscape characteristics, within the practically relevant domain of robotic gripper design.

To achieve this goal, the following tasks are defined:

\begin{itemize}

    \item \textbf{Problem formulation:} Formulation of both gripper configuration problems as multi-objective optimization problems, including the
    definition of decision variables, objective functions, and constraints.
    
    \item \textbf{Optimization framework implementation:} Implementation of a generalized optimization framework capable of handling multi-objective
    problems of varying dimensionality and interfacing with both evaluation pipelines through a unified black-box interface.

    \item \textbf{Evaluation pipeline development:} Development of dedicated black-box evaluation pipelines for each gripper type, capable of
    assessing grasp quality given a gripper configuration and a target workpiece, and of serving as the objective function interface for all
    optimization algorithms.
    
    \item \textbf{Workpiece selection and experimental coverage:} Selection of a diverse set of workpieces spanning a range of geometric properties
    including varying surface curvatures, sizes, and shapes to ensure that algorithm comparisons are not biased towards a specific object geometry.
    Workpiece diversity is treated as a deliberate experimental design choice to expose differences in algorithm robustness and generalization across
    objective landscapes of varying complexity.

    \item \textbf{Algorithm application and comparison:} Application and systematic comparison of a selected set of multi-objective optimization
    algorithms drawn from the available literature \cite{pymoo, olson2025ax, botorch}, evaluated under controlled experimental conditions using
    well-defined performance indicators across both gripper types and workpieces\cite{Audet2018Performance}.

    \item \textbf{Result analysis:} Analysis and interpretation of algorithm performance across both gripper types and all workpieces, with particular
    focus on convergence behaviour, Pareto front quality, and the sensitivity of algorithm rankings to workpiece geometry and objective space
    dimensionality.
\end{itemize}